\subsection{The exact inverse limit of the tail replacements}

\begin{proposition}[A-S09-25: Proved; the natural inverse towers have no admissible cone]
\label{AS0925}
Let $A\in\mathcal A_{\mathbb Z}$, and let
$1\leq N_1\leq N_2\leq\cdots$ tend to infinity.
Use the actual tail functors $T_N$ and the actual skeleton
replacement functor $R$ of A-S09-24 and A-006.
In the ambient category of nonnegatively weight-graded strict
simplicial PD differential algebras over $S=\Gamma_{\mathbb Z}(T)$,
the two inverse systems
\[
 \{T_{N_j}A\}_{j\geq1},\qquad
 \{R T_{N_j}A\}_{j\geq1}
 \tag{S09.125}
\]
have strict inverse limit canonically isomorphic to the
constant coefficient theory $S[0]$.
The transition maps of the first system are those of
(S09.118); those of the second are their images under $R$.
Neither system admits a compatible cone from an admissible
theory.  In particular neither has an inverse limit inside
the admissible category.  These assertions also hold in
the support-restricted and termwise-flat variants.
\end{proposition}

\begin{proof}
Write $V_j=T_{N_j}A$ and $U_j=R T_{N_j}A$.
For $k\geq j$ the transition $V_k\to V_j$ is the
identity below weight $N_j$ and above weight $N_k$;
in intermediate weights it is the coefficient map to $A$.
These actual strict maps satisfy the inverse-system identities
by their componentwise formulas.  Functoriality of $R$
proves the same identities for $U_k\to U_j$.
Every object has its structural coefficient map from $S$,
and all transitions preserve it.  We therefore have actual
compatible cones
\[
                 S\longrightarrow V_j,\qquad
                 S\longrightarrow U_j.
 \tag{S09.126}
\]

Fix a weight $w\geq0$, a form degree $q\geq0$ and a
simplex degree $n\geq0$.  Once $N_j>w$, the defining
formula (S09.116) identifies
\[
 (V_j)_n^{w,q}=
 \begin{cases}S^w,&q=0,\\0,&q>0.\end{cases}
 \tag{S09.127}
\]
The transition maps in this eventual range are identities.
The same equality holds for $U_j$.  Indeed there are no
coefficient-edge fillers of positive weight less than $N_j$
in $V_j$: in such a weight both face maps are the identity
of $S^w$, whereas a filler would have faces $0,T^{[w]}$.
Thus every generator in the free first skeleton $P_{V_j}$
has weight at least $N_j$.  Its matching coskeleton and
the pullback defining $R(V_j)$ have only coefficient
pieces below $N_j$, including weight zero.
This is the explicit low-weight identification used in
(S09.121), and the functorial transition map restricts to
the identity of those coefficient pieces.

Consequently each component of either inverse limit is
canonically $S^w$ in form degree zero and zero in every
positive form degree.  This computes limits in the
weight-graded category: we take componentwise limits and
their graded direct sum, with no completion across weights.
To check all operations, choose $j$ with $N_j$ greater
than every input and output weight of the given operation.
At that stage it is literally the coefficient operation.
This proves that products, divided powers, differentials
and all simplicial operators on the limit are exactly
those of the constant theory $S[0]$.

Here is also its universal property with actual maps.
For any compatible cone of strict maps $f_j:D\to V_j$,
the eventual component (S09.127) defines a unique
weight- and form-preserving map $f:D\to S[0]$.
Compatibility makes this definition independent of $j$.
The preceding operation check proves that $f$ is a strict
simplicial PD cochain map over $S$, and its composites with
(S09.126) are every $f_j$.  Uniqueness follows componentwise
from (S09.127).  The identical proof applies to a cone
$D\to U_j$.  Thus (S09.126) exhibits both strict limits.

Finally no admissible $D$ has a strict coefficient-preserving
simplicial map $f:D\to S[0]$.  Choose
$w\geq\max\{1,\mu_D(0)\}$.
Normalized horizontal row-zero exactness supplies
$a\in D_1^{w,0}$ with
\[
                      d_1a=0,\qquad d_0a=T^{[w]}.
 \tag{S09.128}
\]
Both faces on the constant target $S^w$ are the identity.
Preservation of faces forces $f(a)=0$ from the first
equation, and $f(a)=T^{[w]}$ from the second.
The free coefficient group $S^w=\mathbb Z T^{[w]}$
makes these equations inconsistent.
Hence an admissible compatible cone to either system
cannot exist by the universal property just proved.
The constant theory itself also has normalized horizontal
$H_0=S^w$ in every weight, so it fails the required eventual
horizontal exactness.
Every calculation retained the support and flatness
restrictions, completing the variants.
\end{proof}

This determines the exact outcome of the displayed inverse
systems.  The common-source and terminal-localization gates
of A-S09-24 remain open.
